\section{Conclusion}

We have shown that a passing algorithmic audit can be the proximate cause of a system's subsequent decline: certification's effect on a firm's own investment incentives, distinct from measurement error or threshold miscalibration. Under recurring audits this is a structural choice between a permanently degraded steady state and perpetual cycling, both strictly worse long-run than no safe harbor at all. The degraded steady state is invisible to a regulator observing only pass/fail records, since such a firm never subsequently fails; cycling surfaces as a repeated fail-then-pass pattern, but nothing flags why. The mechanism survives firm heterogeneity and a different functional-form pair; cycling turns out, unexpectedly, to disappear once the threshold is observed imperfectly — but only once noise exceeds an identified threshold, below which it persists and converges to the hard-threshold model's behavior. These predictions are falsifiable, giving regulators a concrete additional design target: safe-harbor strength, audit precision, and threshold choice together.

\paragraph{Limitations.} The results are primarily theoretical: functional forms and parameters are chosen for tractability, not fit to a deployed regime, and Section~\ref{sec:calibration} anchors the threshold against real disclosure statistics without testing the dynamic predictions directly, since the needed firm-level panel does not yet exist. The model assumes myopic optimization throughout; Section~\ref{sec:forwardlooking} shows this is without qualitative loss in Cases (a) and (c), but cycling and the imperfect-observation result remain one-period-ahead simplifications, each holding an auxiliary feature fixed (a representative firm; a stylized detection function) rather than solving the fully general problem, with investment treated as costlessly adjustable throughout.

\paragraph{Future Work.} Three extensions: the joint design problem of choosing $\tau$ and a policy shaping $\rho$ under heterogeneity (Section~\ref{sec:regulator}); relaxing myopic optimization, in the spirit of Che, Kim, and Mierendorff~\cite{che2024datadriven}, for what Proposition~\ref{prop:forwardlooking} does not cover (cycling and imperfect observation, both needing numerical dynamic programming); and testing the kink and cycling predictions directly, awaiting a preserved disclosure panel, constructible if disclosures were archived rather than overwritten.
